\subsection{Unbroken Symmetries}
\label{sec:19.9}

We have seen how the properties of Goldstone bosons and their low energy interactions may be deduced from an assumption that the theory is invariant (or approximately invariant) under a group $G$ spontaneously broken to a subgroup $H$. But in applying these methods to the cases where $G$ is $SU(2)\cdot SU(2)$ or $SU(3)\cdot SU(3)$, we had to take the pattern of symmetry breaking, of $SU(2)\cdot SU(2)$ or $SU(3)\cdot SU(3)$ to their non-chiral $SU(2)$ or $SU(3)$ subgroups, from experiment. In Section 22.5 we shall show that the $SU(3)\cdot SU(3)$ symmetry for massless $u$, $d$, and $s$ quarks must in fact be spontaneously broken in quantum chromodynamics, but it is more difficult to show on the basis of quantum chromodynamics that the $SU(2)\cdot SU(2)$ symmetry with only $u$ and $d$ massless is also spontaneously broken.\footnote{Weingarten~\hyperref[ch19-ref-47]{[47]} has used lattice methods to show that whether or not chiral symmetry is broken, the lightest particle in quantum chromodynamics with massless $u$ and $d$ quarks must have the quantum numbers of the pion. As we will see in Section 22.5, the existence of anomalies due to fermion loops in quantum chromodynamics together with the assumption of quark trapping requires that some hadron be massless, so it follows that the pion is massless, which strongly suggests that chiral symmetry is spontaneously broken.} On the other hand, there is an intuitive argument that their non-chiral $SU(2)$ or $SU(3)$ subgroups are \emph{not} broken, based on a conjecture known as the \emph{persistent mass condition}~\hyperref[ch19-ref-48]{[48]}, which states that composite particles will not be massless if the particles of which they are composed are massive. Non-chiral symmetries like isospin conservation are not violated if we give the quarks equal masses, and then if they were spontaneously broken we would have the massless Goldstone bosons formed as composites of massive quarks, in contradiction with the persistent mass condition. In what follows we shall present a proof by Vafa and Witten~\hyperref[ch19-ref-49]{[49]} that in gauge theories like quantum chromodynamics those non-chiral symmetries that are not violated by quark masses can not be spontaneously broken. This result is of more than academic interest; as we shall see in Section 21.4, it is possible that the spontaneous breakdown of electroweak gauge symmetries is described by a `technicolor' theory similar to quantum chromodynamics, and in testing this idea it is important to know what symmetries are left unbroken in this theory.

Consider a gauge theory like quantum chromodynamics, with a number of fermion `flavors' in identical representations of the gauge group. If all fermions have masses, then the theory will be invariant under all unitary global non-chiral transformations on the fermion flavors that commute with the fermion mass matrix. For instance, if $n_1$ fermions are degenerate with a common mass $m_1$, $n_2$ fermions are degenerate with some other common mass $m_2$, and so on, then this global symmetry group is $U(n_1)\cdot U(n_2)\cdot\cdots$. (As a special case, if there are no degeneracies we have a global symmetry under $U(1)\cdot U(1)\cdot\cdots$; an example is the conservation of baryon number, strangeness, etc. in quantum chromodynamics.) \emph{These symmetries cannot be spontaneously broken.}

To prove this, let us consider a general Greens function for $r$ fermion and antifermion fields.\footnote{In the original work of Vafa and Witten~\hyperref[ch19-ref-49]{[49]}, they first proved the absence of symmetry breaking in the case $r=1$ with $x=y$, then observed that the absence of symmetry breaking in this vacuum expectation value did not rule out the possibility of a spontaneous symmetry breakdown occurring in other Greens functions, and so went on to different methods of proof.} In the path-integral formalism this is\footnote{As we saw in Section 15.5, both numerator and denominator are proportional to the infinite volume of the gauge group, which cancels in the ratio \eqref{eq:19.9.1}. The presence of this infinite factor detracts from the rigor of the following arguments, but if we were to remove it by introducing ghosts then some of the steps below that depend on the positivity of the action would raise difficulties. One way to deal with this problem is to replace the spacetime continuum with a finite lattice of points, in which case the gauge group has a finite volume, and no gauge fixing or ghosts are needed.}
\begin{align}
&\bra{\Omega}T\left\{
\Psi_{u_1k_1}(x_1)\cdots\Psi_{u_rk_r}(x_r)
\Psi^\dagger_{v_1l_1}(y_1)\cdots\Psi^\dagger_{v_rl_r}(y_r)
\right\}\ket{\Omega}
\notag\\
&\quad=\frac{1}{Z}\int[dA][d\psi][d\psi^\dagger]\,
\psi_{u_1k_1}(x_1)\cdots\psi_{u_rk_r}(x_r)
\psi^\dagger_{v_1l_1}(y_1)\cdots\psi^\dagger_{v_rl_r}(y_r)
\notag\\
&\qquad\quad{}\cdot
\exp\left(iI_{\mathrm{gauge}}[A]
+iI_{\mathrm{Dirac}}[\psi,\psi^\dagger;A]\right),
\tag{19.9.1}\label{eq:19.9.1}
\end{align}
where $k_1\ldots k_r$ and $l_1\ldots l_r$ are Dirac spin indices, $u_1\ldots u_r$ and $v_1\ldots v_r$ are flavor indices, $I_{\mathrm{gauge}}[A]$ is the action for a pure gauge theory, $I_{\mathrm{Dirac}}[\psi,\psi^\dagger;A]$ is the action for the Dirac fields in the presence of a gauge field $A^\mu_\alpha(x)$, and $Z$ is the vacuum--vacuum amplitude
\begin{equation}
Z\equiv\int[dA][d\psi][d\psi^\dagger]\,
\exp\left(iI_{\mathrm{gauge}}[A]
+iI_{\mathrm{Dirac}}[\psi,\psi^\dagger;A]\right).
\tag{19.9.2}\label{eq:19.9.2}
\end{equation}

It will be necessary here to work in a Euclidean spacetime, with $x^4=x_4=ix^0$, $y^4=y_4=iy^0$, and $A^4_\alpha=A_{4\alpha}=iA_{0\alpha}$, all real. (See Appendix A of Chapter 23.) In this case, the Dirac action is
\begin{equation}
I_{\mathrm{Dirac}}[\psi,\psi^\dagger;A]
=i\int d^3x\int dx^4\,
\psi^\dagger[\sl{D}+M]\psi,
\tag{19.9.3}\label{eq:19.9.3}
\end{equation}
where $M$ is the fermion mass matrix, and $\sl{D}$ is the Euclidean gauge-covariant derivative contracted with the Euclidean Dirac matrices
\begin{equation}
\sl{D}
=\sum_{i=1}^4(\partial_i-it^\alpha A_{i\alpha})\gamma_i,
\tag{19.9.4}\label{eq:19.9.4}
\end{equation}
where as usual $\gamma_4=i\gamma^0$. Because the action is quadratic in fermion fields, we can explicitly perform the integral over these fields
\begin{align}
&\bra{\Omega}T\left\{
\Psi_{u_1k_1}(x_1)\cdots\Psi_{u_rk_r}(x_r)
\Psi^\dagger_{v_1l_1}(y_1)\cdots\Psi^\dagger_{v_rl_r}(y_r)
\right\}\ket{\Omega}
\notag\\
&\quad=\frac{1}{Z}\int[dA]\,
\operatorname{Det}(\sl{D}+M)
\exp\left(iI_{\mathrm{gauge}}[A]\right)
\notag\\
&\qquad\quad{}\cdot
\left[
(\sl{D}+M)^{-1}_{x_1u_1k_1,y_1v_1l_1}
\cdots
(\sl{D}+M)^{-1}_{x_ru_rk_r,y_rv_rl_r}
\mathbin{\pm}\text{ permutations}
\right],
\tag{19.9.5}\label{eq:19.9.5}
\end{align}
where `$\mathbin{\pm}$ permutations' indicates that we must sum over all $r!$ permutations of the $\psi$ fields, with a minus sign for odd permutations, and
\begin{equation}
Z=\int[dA]\,
\operatorname{Det}(\sl{D}+M)
\exp\left(iI_{\mathrm{gauge}}[A]\right).
\tag{19.9.6}\label{eq:19.9.6}
\end{equation}
The expression \eqref{eq:19.9.5} is manifestly invariant under any unitary transformation on flavor indices that commutes with the mass matrix $M$. It does not matter what non-perturbative effects are produced by the functional integral over gauge fields and ghosts; these fields are inert under the symmetries in question here, so that the remaining functional integral in Eq.~\eqref{eq:19.9.5} cannot break these symmetries.\footnote{It is more difficult to show that $P$, $C$, and $T$ are not spontaneously broken in quantum chromodynamics with massive quarks, because these symmetries act non-trivially on the gauge fields. Vafa and Witten~\hyperref[ch19-ref-50]{[50]} have applied the methods of Ref.~\hyperref[ch19-ref-49]{[49]} to show that $P$ is not spontaneously broken in quantum chromodynamics.} But for this argument to be convincing, we must show that the expression \eqref{eq:19.9.5} is well defined.

This is not merely an academic question of mathematical rigor. As we saw in Section 19.1, the sign that a symmetry is not spontaneously broken is not just that the ground state is invariant under the symmetry transformations, but also that the symmetry of the ground state is stable under small perturbations. If we break the symmetries of $M$ by adding a small perturbation $\delta M$, and if the expression \eqref{eq:19.9.5} becomes singular as $\delta M\to0$, then factors of $\delta M$ in symmetry-breaking matrix elements may be cancelled by the singularities in the matrix elements for $\delta M=0$. Precisely this happens for chiral symmetries, which arise in a symmetry limit where some of the eigenvalues of $M$ vanish. In this case, as we approach the symmetry limit of zero mass, the factors of mass in the numerators of symmetry-breaking expectation values are cancelled by factors of mass in the denominators of the propagators $(\sl{D}+M)^{-1}$ wherever $\sl{D}$ has zero eigenvalues. For this reason, the arguments here will apply only to the non-chiral symmetries of theories with mass matrices $M$ whose eigenvalues are all non-zero.

To set a bound on the matrix element \eqref{eq:19.9.5} in this case, note first that the differential operator \eqref{eq:19.9.4} is here antihermitian, so as long as the hermitian matrix $M$ has no null eigenvalues, $\sl{D}+M$ has a well-defined inverse. It is still necessary to show that the remaining integration over the gauge fields in Eq.~\eqref{eq:19.9.5}, including that in $Z$, does not make this expression singular when $M$ satisfies the conditions for a symmetry, for instance that $n$ of its eigenvalues are equal for a $U(n)$ symmetry. As we shall see, this will insure that when a fermion mass matrix $M$ that is invariant under some global symmetry transformation is perturbed by adding a small term $\delta M$ that breaks this symmetry, the change in the expectation value \eqref{eq:19.9.4} under this symmetry transformation vanishes in the limit $\delta M\to0$.

It is difficult to show that the coordinate-space Greens functions Eq.~\eqref{eq:19.9.5} are non-singular, so let us consider instead the matrix element for smeared fields
\begin{equation}
\Psi_u[f]\equiv\int d^4x\,f^{k\dagger}(x)\Psi_{uk}(x),
\tag{19.9.7}\label{eq:19.9.7}
\end{equation}
where $f^k(x)$ are arbitrary smooth square-integrable functions. In these terms, Eq.~\eqref{eq:19.9.5} reads
\begin{align}
&\bra{\Omega}T\left\{
\Psi_{u_1}[f_1]\cdots\Psi_{u_r}[f_r]
\Psi^\dagger_{v_1}[g_1]\cdots\Psi^\dagger_{v_r}[g_r]
\right\}\ket{\Omega}
\notag\\
&\quad=\frac{1}{Z}\int[dA]\,
\operatorname{Det}(\sl{D}+M)
\exp\left(iI_{\mathrm{gauge}}[A]\right)
\notag\\
&\qquad\quad{}\cdot
\left[
(\sl{D}+M)^{-1}_{f_1u_1,g_1v_1}
\cdots
(\sl{D}+M)^{-1}_{f_ru_r,g_rv_r}
\mathbin{\pm}\text{ permutations}
\right],
\tag{19.9.8}\label{eq:19.9.8}
\end{align}
where
\begin{align}
(\sl{D}+M)^{-1}_{fu,gv}
\equiv{}&
\int d^4x\int d^4y\,
f^{k\dagger}(x)
(\sl{D}+M)^{-1}_{xuk,yvl}
g^l(y).
\tag{19.9.9}\label{eq:19.9.9}
\end{align}
In this basis, the fermion propagators for a given gauge field are not only well defined, but bounded uniformly in $A^\mu_\alpha(x)$:\footnote{To see this, we may expand in the eigenvectors of $M$
\[
M_{uv}c_v^a=m_ac_u^a,
\qquad
c_u^{a*}c_u^b=\delta_{ab},
\]
and, adapting a trick of Vafa and Witten~\hyperref[ch19-ref-49]{[49]}, write Eq.~\eqref{eq:19.9.9} as
\begin{align*}
(\sl{D}+M)^{-1}_{fu,gv}
={}&\sum_a c_u^{a*}c_v^a
\int d^4x\int d^4y\,
f^{k\dagger}(x)
(\sl{D}+m_a)^{-1}_{xk,yl}g^l(y)
\\
={}&\sum_a\mathbin{\pm}c_u^{a*}c_v^a
\int_0^\infty\exp(-|m_a|\tau)
\int d^4x\int d^4y\,
f^{k\dagger}(x)
\bigl(\exp(\mp\tau\sl{D})\bigr)_{xk,yl}
g^l(y),
\end{align*}
where $\pm$ is the sign of $m_a$. The matrix $c_u^{a*}c_v^a$ and the operator $\bigl(\exp(\mp\tau\sl{D})\bigr)_{xk,yl}$ are both unitary, so
\[
\left|
c_u^{a*}c_v^a
\int d^4x\int d^4y\,
f^{k\dagger}(x)
\bigl(\exp(\mp\tau\sl{D})\bigr)_{xk,yl}
g^l(y)
\right|^2
\leq
\int d^4x\,|f(x)|^2
\int d^4y\,|g(y)|^2
=1.
\]
Using this inside the integral yields Eq.~\eqref{eq:19.9.10}.}
\begin{equation}
\left|(\sl{D}+M)^{-1}_{fu,gv}\right|
\leq\sum_a\frac{1}{|m_a|},
\tag{19.9.10}\label{eq:19.9.10}
\end{equation}
where $m_a$ are the eigenvalues of $M$, and for convenience we are now normalizing the functions $f_i(x)$ and $g_j(y)$ so that
\begin{equation}
\int f_i^\dagger(x)f_i(x)\,d^4x
=\int g_j^\dagger(y)g_j(y)\,d^4y
=1.
\tag{19.9.11}\label{eq:19.9.11}
\end{equation}
(The summation convention is suspended here.) Furthermore, the weight function for the average over gauge fields is positive:\footnote{In the fermion determinant, we use the fact that $\operatorname{Det}(\sl{D}+M)=\operatorname{Det}\gamma_5(\sl{D}+M)\gamma_5=\operatorname{Det}(-\sl{D}+M)$.}
\begin{align}
&\operatorname{Det}(\sl{D}+M)
\exp\left(iI_{\mathrm{gauge}}[A]\right)
\notag\\
&\quad=
\exp\left(
-\frac{1}{4}\int d^3x\int dx^4
\sum_{i=1}^4\sum_{j=1}^4F_{ij}^2
\right)
\cdot
\sqrt{\operatorname{Det}\left[(i\sl{D})^2+M^2\right]}.
\tag{19.9.12}\label{eq:19.9.12}
\end{align}
With a positive weight function, the average of any function is bounded by the bound of that function. Using the bound \eqref{eq:19.9.10} in Eq.~\eqref{eq:19.9.8}, the matrix element is bounded by
\begin{align}
\left|
\bra{\Omega}T\left\{
\Psi_{u_1}[f_1]\cdots\Psi_{u_r}[f_r]
\Psi^\dagger_{v_1}[g_1]\cdots\Psi^\dagger_{v_r}[g_r]
\right\}\ket{\Omega}
\right|
\leq r!\left[\sum_a\frac{1}{|m_a|}\right]^r.
\tag{19.9.13}\label{eq:19.9.13}
\end{align}
Thus there are no singularities as the symmetry-breaking terms $\delta M$ in $M$ go to zero that could cancel factors of $\delta M$. To see this more explicitly, we can use the same methods to show that if we perturb $M$ by a small symmetry-breaking term $\delta M$, then the term in Eq.~\eqref{eq:19.9.8} of first order in $\delta M$ is bounded by
\begin{align}
&\left|
\delta\bra{\Omega}T\left\{
\Psi_{u_1}[f_1]\cdots\Psi_{u_r}[f_r]
\Psi^\dagger_{v_1}[g_1]\cdots\Psi^\dagger_{v_r}[g_r]
\right\}\ket{\Omega}
\right|
\notag\\
&\qquad\leq
r\,r!\sum_{a,b}
\frac{|(\delta M)_{ab}|}{|m_am_b|}
\left[\sum_a\frac{1}{|m_a|}\right]^{r-1},
\tag{19.9.14}\label{eq:19.9.14}
\end{align}
so this vanishes when $\delta M\to0$.

In the real world none of the quark masses are zero or degenerate, so it follows immediately from the above arguments that the $U(1)$ symmetries like conservation of baryon number, strangeness, etc., are neither spontaneously nor intrinsically broken. The other strong-interaction symmetries like isospin or $SU(3)$ are more problematic. These symmetries would be completely unbroken if the $u$ and $d$ quarks, or $u$, $d$, and $s$ quarks, had equal non-zero masses, but as we saw in Section 19.7, these masses are not at all degenerate; the isospin and $SU(3)$ symmetries arise because the quark masses are small, not equal. These symmetries do remain unbroken if we give two or three quarks equal masses and let the masses become arbitrarily small, so isospin or $SU(3)$ will be good symmetries in any process that is not sensitive to the small quark masses. But not all processes are insensitive to these masses, because for two or three massless quarks, the pion or the pseudoscalar meson octet to which it belongs becomes massless. This is not a problem for isospin, because as we saw in Section 19.5, the pion triplet remains degenerate to first order in quark masses even though $m_u\ne m_d$. For $SU(3)$ the quark mass differences produce first-order mass differences among the pion, kaon, and eta, so that processes dominated by one-meson poles \emph{can} show large departures from $SU(3)$.
